\section{Introduction}

Interactive responsiveness is a tail-latency problem~\cite{dean2013tail}. A user does not notice the average time to handle a keystroke or render a frame; they notice the slow ones. On a desktop, the slow ones are often caused by work the user did not start: a sync client hashing files, a search indexer, an update, a compiler. Operating systems expose resource \emph{utilisation} (Task Manager on Windows), but utilisation is a poor proxy for harm. Memory-bandwidth contention, for example, can make an interactive task several times slower while half the cores sit idle. It is invisible to every utilisation counter Windows provides to user mode.

Linux addressed the measurement side with Pressure Stall Information (PSI)~\cite{weiner2018psi}, which reports the share of time tasks are stalled on CPU, memory, or I/O. Datacentre systems such as Heracles, PARTIES, and Caladan~\cite{lo2015heracles,chen2019parties,fried2020caladan} address the control side with feedback loops over hardware partitioning, assuming privileged control and instrumented services. A Windows desktop has neither: no PSI equivalent, no cache or bandwidth partitioning exposed to applications, no administrator rights for ordinary software, and a foreground that cannot be instrumented.

This paper asks what can be done within those constraints. Our system, SmartSwap, contributes:
\begin{enumerate}
  \item \textbf{A user-mode stall index for Windows (WSI).} A low-duty canary times its own scheduling delay, compute, memory streaming, and unbuffered storage reads. From these timings we derive the share of wall time an interactive task would lose, with no privileges, drivers, ETW sessions, or hardware counters. We validate it against a separate, un-instrumented foreground: $\rho=0.96$, AUC $\geq 0.965$ (\S\ref{sec:results}).
  \item \textbf{Specificity-ordered differential attribution.} We show that a naive ``largest slowdown wins'' rule misattributes memory-bandwidth contention to storage. Ordering signals by specificity (CPU, then memory, then I/O) names the resource correctly in 86--99.6\,\% of stalled periods.
  \item \textbf{A least-cost, doubly verified control ladder.} Scheduler hints come first; an AIMD-governed CPU-rate cap is used only if needed. Each action is checked by kernel read-back and by measured benefit, and is rolled back otherwise. A load-aware release rule prevents oscillation. An ablation shows the ladder is essential (up to $4.8\times$ worse p95 without it).
  \item \textbf{A cost-aware, randomised-block evaluation.} Every policy is measured on foreground tail latency \emph{and} on background work lost. Comparisons are paired and Holm-adjusted, the control is genuinely free of interference, and validity checks confirm the injected load actually interferes. We include the strongest static baseline and report where it equals or beats our controller.
  \item \textbf{A post-mortem of a flawed evaluation.} Our first evaluation reported significant improvements produced by six silent defects, including a workload that never ran. We describe each defect and the guard that now prevents it (\S\ref{sec:lessons}).
\end{enumerate}
\S\ref{sec:design} describes the design and \S\ref{sec:impl} the implementation. \S\ref{sec:method} gives the methodology and \S\ref{sec:results} the results. \S\ref{sec:discussion} discusses limitations.
